\chapter{Implementation}

\section{Backend implementation}

The backend is organized into modules by business domain (users, machines, work orders, intervention orders, alerts, scheduling, Machine Learning). Each module exposes its API routes, validation schemas, and business services separately, which facilitates maintenance and future changes. Role-based access control is applied centrally on every route. The four workflows below (maintenance planning, intervention requests, work orders, and inventory reservations) illustrate this modular organization at the business-logic level.

\subsection{Maintenance planning workflow}

An Administrator opens a planning window in \emph{draft} status; the Head Technician then breaks it down into per-machine diagnostic or correction tasks and submits it for approval. Only once an Administrator approves the planning are the assigned technicians notified, by both an in-app alert and an email.

\begin{figure}[H]
\centering
\includegraphics[width=\textwidth]{figures/planning.png}
\caption{Maintenance planning workflow}
\label{fig:planning-workflow}
\end{figure}

\begin{figure}[H]
\centering
\includegraphics[width=\textwidth]{figures/ss_planning.png}
\caption{Planning screen (Administrator)}
\label{fig:ss-planning}
\end{figure}

\begin{figure}[H]
\centering
\includegraphics[width=\textwidth]{figures/ss_planning_2.png}
\caption{Planning screen: per-machine task breakdown}
\label{fig:ss-planning-2}
\end{figure}

\subsection{Intervention request workflow}

A field technician's request passes through a single Head Technician validation gate with two distinct outcomes: approval requires an associated machine and automatically creates a linked work order, while rejection requires a reason and releases any parts reservation already made. Closing the loop later, once the work is done, still requires a root-cause diagnosis before the intervention can be marked validated.

\begin{figure}[H]
\centering
\includegraphics[width=\textwidth]{figures/intervention.png}
\caption{Intervention request workflow}
\label{fig:intervention-workflow}
\end{figure}

\begin{figure}[H]
\centering
\includegraphics[width=\textwidth]{figures/ss_intervention.png}
\caption{Intervention orders screen}
\label{fig:ss-intervention}
\end{figure}

\subsection{Work order lifecycle}

A work order moves through a sequential ten-state lifecycle. It can be rejected at the approval step, releasing any linked stock reservation, and it can only be closed by an Administrator once it has reached the \emph{validated} state.

\begin{figure}[H]
\centering
\includegraphics[width=\textwidth]{figures/workorder.png}
\caption{Work order lifecycle}
\label{fig:wo-lifecycle}
\end{figure}

\begin{figure}[H]
\centering
\includegraphics[width=\textwidth]{figures/ss_workorder.png}
\caption{Work order management screen}
\label{fig:ss-workorder}
\end{figure}

\begin{figure}[H]
\centering
\includegraphics[width=\textwidth]{figures/ss_wo_closer.png}
\caption{Work order close-out form (Technician)}
\label{fig:ss-wo-closer}
\end{figure}

\subsection{Inventory reservation system}

A single service owns every mutation of stock quantities for the intervention/work-order workflow. Four independent triggers (approval, rejection, a daily scheduled task, and technician completion) each invoke a dedicated process, all reading and writing the same reservation records; approval handlers never touch stock tables directly.

\begin{figure}[H]
\centering
\includegraphics[width=\textwidth]{figures/inventory.png}
\caption{Inventory reservation data flow}
\label{fig:inventory-flow}
\end{figure}

\begin{figure}[H]
\centering
\includegraphics[width=\textwidth]{figures/ss_stock.png}
\caption{Spare parts inventory screen}
\label{fig:ss-stock}
\end{figure}

\subsection{Summary}

The four core workflows (planning, intervention requests, work orders, inventory reservations) are each owned by a single dedicated service, keeping stock mutations, approval logic, and state transitions from leaking into unrelated modules. This domain-by-domain split is what lets a rule change in one workflow (for example, how stock is reserved) ship without touching the others.

\section{Frontend implementation}

The frontend offers three distinct application spaces, each tailored to the needs of the corresponding role: an Administrator space for global management, a Head Technician space for oversight and scheduling, and a Technician space for carrying out interventions in the field. A consistent design system (color palette, reusable components) ensures a uniform experience across the three spaces.

\begin{figure}[H]
\centering
\includegraphics[width=\textwidth]{figures/ss_login.png}
\caption{Login screen}
\label{fig:ss-login}
\end{figure}

\begin{figure}[H]
\centering
\includegraphics[width=\textwidth]{figures/ss_dashboard_admin.png}
\caption{Administrator dashboard}
\label{fig:dashboard-admin}
\end{figure}

\begin{figure}[H]
\centering
\includegraphics[width=\textwidth]{figures/ss_approval_users.png}
\caption{User approval screen (Administrator)}
\label{fig:ss-approval-users}
\end{figure}

\begin{figure}[H]
\centering
\includegraphics[width=\textwidth]{figures/ss_approval_itv.png}
\caption{Intervention approval screen (Head Technician/administrator)}
\label{fig:ss-approval-itv}
\end{figure}

\begin{figure}[H]
\centering
\includegraphics[width=\textwidth]{figures/ss_dashboard_cheftech.png}
\caption{Head Technician dashboard}
\label{fig:dashboard-cheftech}
\end{figure}

\begin{figure}[H]
\centering
\includegraphics[width=\textwidth]{figures/ss_technicien.png}
\caption{Technician dashboard}
\label{fig:ss-technicien}
\end{figure}

\begin{figure}[H]
\centering
\includegraphics[width=\textwidth]{figures/ss_machine_detail.png}
\caption{Machine detail and health screen}
\label{fig:ss-machine-detail}
\end{figure}

\subsection{Summary}

Three role-scoped spaces (Administrator, Head Technician, Technician) share one design system, from the shared login screen through to each role's dashboard, approval queues, and machine-detail views. Keeping the UI role-aware at the routing level, not just visually, is what enforces the access boundaries defined in Chapter~2 on the client side, on top of the server-side checks.

\section{Machine Learning pipeline implementation}

Each predictive model (P1 to P6) was trained and then integrated into the model-serving microservice, with lazy, cached loading to limit response times. An explainability layer (``Why?'') was added across the board: every prediction, alert, or recommendation can be expanded into a plain-language explanation, free of Machine Learning terminology, so it remains understandable to field technicians.

\begin{figure}[H]
\centering
\includegraphics[width=\textwidth]{figures/ss_ml_dashboard.png}
\caption{ML predictions screen}
\label{fig:ss-ml-dashboard}
\end{figure}

\begin{figure}[H]
\centering
\includegraphics[width=\textwidth]{figures/ss_alerts.png}
\caption{Predictive alerts screen}
\label{fig:ss-alerts}
\end{figure}

The P7 spare parts coordination module was built in six progressive steps: a demand forecasting engine, stock shortage alert routing, role-based interface adaptation, guided order preparation (never automatic), an operational readiness score with an event timeline, and a feedback loop comparing predictions against parts actually consumed.

\begin{figure}[H]
\centering
\includegraphics[width=\textwidth]{figures/ss_procurement.png}
\caption{Guided procurement screen}
\label{fig:ss-procurement}
\end{figure}

\begin{figure}[H]
\centering
\includegraphics[width=\textwidth]{figures/ss_explicabilite.png}
\caption{Prediction explainability panel}
\label{fig:explicabilite}
\end{figure}

\subsection{Summary}

P1--P6 are served through a lazily loaded, cached model registry, with every prediction, alert, and recommendation expandable into a plain-language ``Why?'' explanation so field technicians never have to interpret raw model output. P7's guided procurement flow is the clearest expression of the platform's human-in-the-loop principle: the system forecasts and proposes, but never places an order on its own.

\section{RAG service and conversational bridge implementation}

The RAG service supports importing technical documents (manuals, procedures) with version management and cascading deletion. The bridge between the RAG service and the ML microservice was designed to never let a conversation fail: if the ML microservice is unavailable, a fallback estimate based on simple rules is used instead, transparently to the user.

\begin{figure}[H]
\centering
\includegraphics[width=\textwidth]{figures/ss_ragdocs.png}
\caption{Knowledge base document management screen}
\label{fig:ss-ragdocs}
\end{figure}

\begin{figure}[H]
\centering
\includegraphics[width=\textwidth]{figures/ss_chat.png}
\caption{Conversational assistant screen}
\label{fig:ss-chat}
\end{figure}

\subsection{Summary}

Document management (versioning, cascading deletion) and the ML bridge were built together so the two failure modes stay independent: a bad document upload never breaks predictions, and a down ML microservice never breaks the chat, since the bridge falls back to a rule-based estimate instead of failing the conversation.

\section{Pilot rollout plan (PDCA)}

Deploying six new predictive models directly to every machine in the plant, with no feedback loop, would risk both poor predictions and loss of user trust. The rollout was therefore planned as a Plan-Do-Check-Act (PDCA) cycle: a controlled, measurable pilot covering the Plan and Do phases below, before any decision to generalize.

\paragraph{Plan.} The pilot targeted a single functional zone (surface-mount technology, a production bottleneck with high-frequency, reliable sensor data) for eight weeks, split into three weeks of \emph{shadow mode} (predictions computed and logged, but not shown to users) followed by five weeks of \emph{active mode}, where the dashboard is revealed to two maintenance managers and four senior technicians. Four quantitative success criteria were fixed in advance: at least 15\% reduction in unexpected downtime against a three-month historical baseline; at least 80\% of machines flagged \emph{critical} by the failure-probability model confirmed as a real issue on inspection; at least 75\% agreement between the predicted and the technician-confirmed failure type; and no more than 250\,ms of additional backend latency introduced by the inference calls.

\paragraph{Do.} During shadow mode, every prediction was written to a dedicated log table and compared, at the end of the three weeks, against the intervention orders that managers actually created: a first check of whether the models would have caught what really happened. During active mode, managers were required to open an intervention order for every machine flagged \emph{critical}, and technicians were required, when closing a work order, to record the actual failure type observed. This technician-confirmed label capture is the direct precursor of the production feedback loop described in Chapter~5: real, technician-confirmed outcomes compared against past predictions, closing the loop between the field and the model.

\begin{figure}[H]
\centering
\includegraphics[width=\textwidth]{figures/ss_pdca.png}
\caption{Pilot rollout dashboard (PDCA tracking)}
\label{fig:ss-pdca}
\end{figure}

\section{Production deployment}

Development up to this point ran entirely on Docker Compose. To validate the platform under conditions closer to a real operating environment (container orchestration, a managed database, centralized secret management, and public network access), the full stack was additionally deployed to a managed Kubernetes cluster on Microsoft Azure, provisioned as Infrastructure as Code. Figure~\ref{fig:deployment-pipeline} summarizes the resulting pipeline, from infrastructure provisioning down to the public web address; each step is explained below in plain terms.

\begin{figure}[H]
\centering
\includegraphics[width=0.85\textwidth]{figures/deployment_pipeline.png}
\caption{Production deployment pipeline, from infrastructure to public access}
\label{fig:deployment-pipeline}
\end{figure}

\paragraph{Terraform.} Every cloud resource this platform runs on is declared in Terraform's HCL syntax rather than clicked together in the Azure Portal, so the whole infrastructure can be rebuilt from the same text files already under version control: no undocumented manual step, and every change goes through the same plan/apply review a code change would.

\paragraph{Azure resources.} Three resources come out of that Terraform configuration: the \textbf{Azure Kubernetes Service (AKS)} cluster the application runs on, the managed PostgreSQL Flexible Server behind it, and a Key Vault holding the database password and other secrets away from both the application code and the Helm values files.

\paragraph{Helm chart.} The seven application components, backend, frontend, ML microservice, RAG service, Celery worker, Celery beat, RabbitMQ, are packaged into a single Helm chart, so a full environment comes up from one install command instead of applying seven sets of Kubernetes manifests by hand and keeping them in sync manually.

\paragraph{Staging vs.\ production.} The same package is installed twice, into two separate, isolated environments. \emph{eam-staging} is loaded with realistic demonstration data and used for testing and screenshots. \emph{eam-prod} starts instead from a clean, empty database, with no demonstration data mixed in, so it reflects a genuine first-run deployment.

\paragraph{Public access.} Public exposure is deliberately scoped to production only. Inside that cluster, a lightweight tunnel agent keeps an outbound-only connection alive to a relay service; that outbound link is what lets an outside visitor reach the app over HTTPS at a fixed address, without ever opening an inbound port on the cluster itself.

\begin{figure}[H]
\centering
\includegraphics[width=\textwidth]{figures/aks_deployment.png}
\caption{Production Kubernetes cluster (Azure Portal)}
\label{fig:aks-deployment}
\end{figure}

The cluster, the managed PostgreSQL server (with the pgvector extension enabled for the RAG service), and a Key Vault for secret storage were all provisioned with Terraform, with remote state kept on \textbf{HashiCorp Cloud Platform (HCP)} Terraform for visibility into every planned and applied change.

\begin{figure}[H]
\centering
\includegraphics[width=\textwidth]{figures/terraform.png}
\caption{Terraform-managed infrastructure (HCP Terraform workspace)}
\label{fig:terraform-iac}
\end{figure}

That same bundle, backend, frontend, ML microservice, RAG service, the two Celery processes, and RabbitMQ, together with MinIO for object storage, was deployed twice from one chart with different value files: once into a staging namespace seeded with realistic demonstration data, and once into a separate, clean production namespace with its own database and no demonstration data. Pods authenticate to Key Vault through Azure Workload Identity, via the Secrets Store CSI Driver, so no credential is ever stored in the chart's values or in a Kubernetes Secret checked into version control.

\begin{figure}[H]
\centering
\includegraphics[width=\textwidth]{figures/helm.png}
\caption{Helm releases across the cluster: staging and production side by side}
\label{fig:helm-dashboard}
\end{figure}

Public HTTPS access is provided by a lightweight tunnel agent running inside the cluster, which exposes the production environment on a stable public address. The production environment is reachable publicly over HTTPS through this route.

\section{Technical difficulties encountered}

Several difficulties were encountered during implementation. Synchronizing trained models between the research environment (notebooks) and the production environment (microservice) required setting up a strict procedure for copying and verifying the number of features expected by each model, in order to avoid loading a stale model incompatible with the data format sent in production.

Partial component availability was also a point of attention: since certain deep learning libraries were not installed in the production environment, the anomaly module's autoencoder was made optional, with automatic renormalization of the detection ensemble's weights when this component is absent. Likewise, the ML-RAG bridge was designed to fall back to a simple rule-based estimate if the ML microservice is unavailable, ensuring a response is always provided to the user.

Finally, the consistency of large-scale test data (several thousand simulated records) required systematic verification of the representativeness of the generated failure scenarios, so that evaluation metrics would reflect realistic cases.
